%% ============================================================
%%  A4 -- Phase 5 Signal Research: BTCUSDT Perpetual Futures
%%  Compile with: pdflatex A4.tex  (twice for refs)
%% ============================================================
\documentclass[12pt, a4paper]{article}

\usepackage{amsmath, amssymb, amsthm, mathtools}
\usepackage{booktabs}
\usepackage{hyperref}
\usepackage{geometry}
\usepackage{xcolor}
\usepackage{array}
\usepackage{multirow}
\usepackage{enumitem}
\usepackage{caption}
\usepackage{float}
\usepackage{parskip}
\usepackage{bm}

\geometry{margin=2.5cm}

\hypersetup{
  colorlinks=true,
  linkcolor=blue!60!black,
  urlcolor=blue!60!black
}

\theoremstyle{plain}
\newtheorem{theorem}{Theorem}[section]
\newtheorem{proposition}{Proposition}[section]
\newtheorem{corollary}{Corollary}[section]
\theoremstyle{definition}
\newtheorem{definition}{Definition}[section]
\newtheorem{remark}{Remark}[section]

\newcommand{\IC}{\mathrm{IC}}
\newcommand{\E}{\mathbb{E}}
\newcommand{\Var}{\mathrm{Var}}
\newcommand{\Cov}{\mathrm{Cov}}
\newcommand{\Cor}{\mathrm{Cor}}
\newcommand{\DVOL}{\mathrm{DVOL}}
\newcommand{\bp}{\,\mathrm{bp}}
\newcommand{\skew}{\widehat{\gamma}}

\title{\textbf{Phase 5 Signal Research --- BTCUSDT Perpetual Futures}\\[0.4em]
\large Realized Skewness, Pre-Settlement Flow, and Cross-Exchange Funding Divergence\\[0.3em]
\large A Complete Treatment from First Principles}
\author{}
\date{May 2026}

\begin{document}
\maketitle
\tableofcontents
\newpage

%% ============================================================
\section{Introduction}
%% ============================================================

This document records Phase 5 of the BTCUSDT perpetual futures signal
research programme.  It is a self-contained technical record of the
design, derivation, and empirical testing of three candidate signal
families, each motivated by genuinely new economic mechanisms.

The core result of Phase 5 is \textbf{negative}: no standalone
higher-frequency signal was validated.  Specifically:

\begin{enumerate}
  \item \textbf{S1a (Realized Skewness) passes the IC screen but is
        killed in the deep dive.}  The 95\% bootstrap confidence interval
        crosses zero (lower bound $= -0.008$), the incremental IC after
        controlling for N3$z$ is not significant ($p = 0.234$), and the
        signal loses money ($\text{Sharpe} = -2.26$) on days when N3 is
        not simultaneously active.  S1a's positive IC is an artifact of
        its overlap with the N3 regime.

  \item \textbf{S1b (Skewness velocity), S2a/S2b (Pre-Settlement Flow),
        and S3a (Cross-Exchange Funding Divergence) are all killed at the
        IC screen.}  Each fails at least one of the pre-committed gates.

  \item \textbf{S3b (OKX--Binance divergence) is untestable.}  The OKX
        public API retains only $\approx 90$ days of funding history,
        providing 7\% coverage of the OOS window.
\end{enumerate}

The fundamental finding is that edge in BTCUSDT perpetuals remains
concentrated in high-DVOL fear events (the N3/P3 regime).  Mechanisms
that operate at sub-daily timescales require a different research
infrastructure --- intraday return targets, tick-level data, and explicit
settlement-mechanics modelling --- that is outside the scope of the
current 24h-hold framework.

The validated signal stack at the end of Phase 5 is:
\begin{center}
\begin{tabular}{llll}
\toprule
Signal & Entry rule & OOS Sharpe & Source \\
\midrule
N3 & N3$z > \theta$ AND DVOL $\geq \delta$ & $+2.95$ & A1 \\
P3 & DD regime AND DVOL $\geq \delta$ & $+5.18$ (exclusive) & A2 \\
\bottomrule
\end{tabular}
\end{center}

The report proceeds as follows. Section~\ref{sec:market} introduces
each new mechanism from first principles: the statistical theory of
realized skewness, pre-settlement funding mechanics, and the economics
of cross-exchange funding divergence.  Section~\ref{sec:stats} summarises
the statistical framework.  Sections~\ref{sec:s1}--\ref{sec:s3} analyse
each signal family.  Section~\ref{sec:synthesis} synthesises the finding
on the high-frequency problem.  Section~\ref{sec:conclusion} concludes.

%% ============================================================
\section{Market Background}
\label{sec:market}
%% ============================================================

\subsection{Statistical Moments and Realized Skewness}
\label{sec:skew_theory}

\subsubsection{Statistical Moments}

For a random variable $X$ with distribution $F$, the \emph{$k$-th
central moment} is:
\begin{equation}
  \mu_k = \E\!\left[(X - \E[X])^k\right]
  = \int (x - \mu)^k \,\mathrm{d}F(x)
\end{equation}
The first four moments have names and interpretations:
\begin{itemize}[noitemsep]
  \item $\mu_1 = \E[X]$: the mean (location).
  \item $\mu_2 = \Var(X)$: the variance (spread around the mean).
  \item $\mu_3$: the third central moment (asymmetry).
  \item $\mu_4$: the fourth central moment (tail weight / kurtosis).
\end{itemize}

\begin{definition}[Standardised Skewness]
  The \emph{standardised skewness} (or simply skewness) is the third
  central moment normalised by the cube of the standard deviation:
  \begin{equation}
    \gamma = \frac{\mu_3}{\sigma^3}
           = \frac{\E[(X - \mu)^3]}{\left(\E[(X-\mu)^2]\right)^{3/2}}
    \label{eq:skewness}
  \end{equation}
  $\gamma = 0$ for a symmetric distribution (e.g.\ the Normal).
  $\gamma < 0$ indicates a \emph{left-skewed} (negatively skewed)
  distribution: a long left tail with many small positive returns and
  occasional large negative returns.  $\gamma > 0$ indicates a right-skewed
  distribution.
\end{definition}

\subsubsection{Skewness and Crash Risk}

Harvey \& Siddique (2000) show that conditional skewness is a priced
risk factor in equity markets: assets with negative coskewness with the
market --- meaning they tend to fall most when the market is also falling
and volatility is high --- earn higher expected returns as compensation.

Amaya et al.\ (2015) extend this to the \emph{realised} domain: skewness
computed from high-frequency returns predicts future cross-sectional
equity returns, with negatively skewed assets subsequently
outperforming.  Their interpretation is that negative realised skewness
captures an episode of crash-risk resolution: after a period during which
returns were systematically left-skewed (frequent small gains interrupted
by large losses), the distributional shape tends to normalise.

\subsubsection{Realized Skewness in Crypto Perpetuals}

In BTCUSDT perpetuals, the 1-minute log-return distribution has a
conditional skewness that varies over time.  A sustained period of
negative realised skewness indicates:

\begin{enumerate}
  \item Many small upward ticks (the normal background) interrupted by
        occasional large downward moves.
  \item The large downward moves are consistent with \emph{forced liquidation
        events}: as over-levered longs are liquidated, the engine sells
        into thin markets, creating outsized downward moves relative to
        the background volatility.
  \item Once the liquidation cascade has cleared (over-levered positions
        are gone), the downward skew normalises.
\end{enumerate}

The Signal S1a tests this contrarian hypothesis: negative 5-day realised
skewness predicts positive 24h returns, mediated by the liquidation
exhaustion mechanism.

\subsubsection{Realized vs.\ Implied Skewness}

There is an important distinction between \emph{realized} and
\emph{implied} skewness.  Implied skewness is extracted from the
options market (the slope of the implied volatility smile across strikes).
Realized skewness is computed directly from past price data.  Signal S1
uses realized skewness because:

\begin{itemize}[noitemsep]
  \item High-frequency 1-minute data is available at no cost.
  \item Implied skewness data requires paid access to options tick data.
  \item Realized skewness is a backward-looking measure directly tied
        to the liquidation dynamics in the perpetual contract itself.
\end{itemize}

\subsection{The Funding Settlement Mechanism and Pre-Settlement Flow}
\label{sec:settlement_theory}

\subsubsection{The 8-Hour Settlement Cycle}

As described in A1 (Section 2.1), Binance settles funding every 8 hours
(at 00:00, 08:00, and 16:00 UTC).  At each settlement, the cash transfer
between longs and shorts is:
\begin{equation}
  \text{Payment}_\tau = F_\tau \times \text{Notional}
\end{equation}
where $F_\tau$ is the funding rate at settlement time $\tau$.

\subsubsection{The Pre-Settlement Flow Hypothesis}

In the minutes preceding a settlement, informed market participants face
a specific opportunity:

\begin{enumerate}
  \item If $F_\tau > 0$ (longs will pay), a trader can:
        \begin{enumerate}
          \item Enter short before settlement, receiving funding at $\tau$.
          \item Close the short immediately after settlement.
          \item Net profit: the funding received minus transaction costs
                (entry + exit $= 6\bp$).
        \end{enumerate}
  \item For this to be profitable, $F_\tau > 6\bp$ (which occurs
        occasionally during extreme funding episodes).
\end{enumerate}

At the aggregate market level, pre-settlement shorting creates
\emph{taker sell pressure} in the 30 minutes before settlement.
This supply pressure temporarily depresses the mark price, creating
a predictable post-settlement recovery as the artificial short position
is unwound.

Signal S2a captures the momentum version: pre-settlement selling predicts
downward continuation (the shorts are informed about further decline).
Signal S2b captures the contrarian version: the selling is mechanical
(funding arbitrage) rather than informed, so it reverts after settlement.

\subsubsection{Why Daily-Resolution Testing is Problematic}

The pre-settlement mechanism operates on the timescale of \emph{the
settlement event}: the relevant forward return is the next 2--4 hours
after settlement, not the next 24 hours.  Aggregating the signal to a
daily (24h) evaluation:

\begin{itemize}
  \item The signal fires only at 3 settlement events per day.
  \item The relevant return window (2--4h post-settlement) is entirely
        contained within the 24h measurement horizon, diluted by 20+
        additional hours of returns that are not predicted by this signal.
  \item Three settlement events per day means only $\sim 3$ signal
        observations per day, not one, violating the daily-IC framework.
\end{itemize}

This mismatch between signal resolution and evaluation resolution is
the primary reason S2 is expected to fail the IC screen.
Section~\ref{sec:s2} confirms this expectation empirically.

\subsection{Cross-Exchange Funding Rate Divergence}
\label{sec:funding_div_theory}

\subsubsection{Why Funding Rates Differ Across Exchanges}

Major exchanges (Binance, Bybit, OKX, Deribit) each publish their own
8-hour funding rates for BTCUSDT perpetual futures.  These rates are
determined by each exchange's own perpetual price relative to their own
spot index.  They can differ because:

\begin{enumerate}
  \item \textbf{Speculative demand imbalance:}  Each exchange has a
        distinct user base.  Retail traders on Bybit may be more
        aggressively long than those on Binance during a bull phase,
        driving Bybit funding above Binance funding.
  \item \textbf{Liquidity differences:}  Thinner markets have larger
        bid-ask spreads and higher price impact, allowing larger
        discrepancies from the spot index before arbitrageurs close them.
  \item \textbf{Institutional flow:}  Large institutional orders often
        arrive at a single exchange, creating temporary imbalances that
        persist for the duration of a settlement window.
\end{enumerate}

\subsubsection{The Arbitrage Channel}

Makarov \& Schoar (2020) document persistent price differences across
cryptocurrency exchanges and find that arbitrage capital is insufficient
to close gaps immediately.  The analogous funding rate divergence
argument:

\begin{proposition}[Funding Divergence Arbitrage]
  If Bybit's funding rate $F^{\mathrm{Byb}}_\tau > F^{\mathrm{Bin}}_\tau$,
  an arbitrageur can earn:
  \[
    \text{arb profit} = F^{\mathrm{Byb}}_\tau - F^{\mathrm{Bin}}_\tau
                        - c_{\mathrm{Byb}} - c_{\mathrm{Bin}}
  \]
  by simultaneously going short on Bybit (collecting the higher rate)
  and long on Binance (paying the lower rate).  The position is
  delta-neutral.
\end{proposition}

If this arbitrage is active and large, it creates:
\begin{itemize}[noitemsep]
  \item \textbf{Short pressure on Bybit:} the arb short raises
        Bybit's selling pressure, dragging Bybit funding down.
  \item \textbf{Long pressure on Binance:} the arb long raises
        Binance buying pressure, pushing Binance prices up.
\end{itemize}

The directional prediction for Binance: when Bybit funding exceeds
Binance funding (positive divergence), the arbitrage flow creates
long pressure on Binance, predicting positive Binance returns.

\subsubsection{Limits to Arbitrage in Crypto Funding}

The hypothesis is theoretically sound, but several friction sources
limit its practical strength at the daily frequency:

\begin{itemize}
  \item \textbf{Transaction costs:}  Entering a delta-neutral position
        on two exchanges requires two round-trips ($\approx 12\bp$ total),
        so the funding rate must diverge by $> 12\bp$ to be profitable.
        At the 8-hour horizon, this is a high bar.
  \item \textbf{Cross-exchange settlement misalignment:}  Different
        exchanges may have slightly different settlement times, creating
        basis risk in the arb leg.
  \item \textbf{Liquidity constraints:}  Large arb positions on thinner
        exchanges incur market impact costs.
  \item \textbf{Rapid convergence:}  If arb capital is well-capitalized,
        any divergence closes within minutes (not hours), leaving no
        signal at the daily resolution.
\end{itemize}

Section~\ref{sec:s3} tests empirically whether any residual divergence
persists long enough to predict 24h Binance returns.

%% ============================================================
\section{Statistical Framework}
\label{sec:stats}
%% ============================================================

This section derives the statistical tools used throughout Phase 5 from
first principles.  The treatment is self-contained.

\subsection{Log-Returns as Prediction Targets}

For a price series $\{C_t\}$, the \emph{simple return} over a horizon $h$ is
$\tilde{r}_t^{(h)} = C_{t+h}/C_t - 1$.  The \emph{log-return} is:
\[
  r_t^{(h)} = \log\frac{C_{t+h}}{C_t} = \log C_{t+h} - \log C_t
\]
Log-returns are preferred for three reasons:

\paragraph{Temporal additivity.}
Log-returns over sub-periods add exactly: $r_{t,t+h+k} = r_{t,t+h} + r_{t+h,t+h+k}$.
Simple returns compound multiplicatively, requiring chain multiplication.
This additivity is essential for equity curve construction: the log-equity
curve is the cumulative sum of log-returns.

\paragraph{Symmetry.}
A simple return of $+50\%$ is not cancelled by a subsequent $-50\%$:
$1.5 \times 0.5 = 0.75 \neq 1$.  Log-returns are symmetric: $+\log 1.5$
is exactly cancelled by $-\log 1.5$.  This symmetry makes long and short
positions analytically equivalent, since a short log-return is $-r_t^{(h)}$.

\paragraph{Cost comparability.}
Transaction costs are approximately constant fractions of notional regardless
of direction.  A round-trip cost $c_{\text{RT}}$ is subtracted uniformly from
$r_t^{(h)}$ regardless of its sign or magnitude.  This is exact in log-return
space and approximate in simple return space.

\paragraph{Basis-point convention.}
All PnL, drawdown, and per-trade figures are stated in \emph{basis points}:
$1\bp = 0.01\% = 10^{-4}$.  A log-return of $0.01$ (1\%) equals $100\bp$.

The 24-hour total return including funding payments is:
\begin{equation}
  \pi_t^{(24h)} = \log\frac{C_{t+1440}}{C_t}
                  - \sum_{k:\,\tau_k \in [t, t+1440]} F_{\tau_k}
\end{equation}
where $F_{\tau_k}$ is the 8-hour funding rate at settlement time $\tau_k$,
and the sum covers the three settlements within the hold window.
Maker round-trip cost: $c_{\mathrm{RT}} = 6\bp$ ($3\bp$ entry, $3\bp$ exit).

\subsection{The Spearman Information Coefficient}
\label{sec:IC}

\begin{definition}[Spearman IC]
  Let $\{(s_t, r_t)\}_{t=1}^n$ be a sample of signal-return pairs.
  Define the rank functions:
  \[
    R_t = \rank(s_t) = \left|\{j : s_j \leq s_t\}\right|, \qquad
    Q_t = \rank(r_t) = \left|\{j : r_j \leq r_t\}\right|
  \]
  (ranks from 1 to $n$; midranks for ties).  The Spearman IC is the Pearson
  correlation applied to the rank vectors:
  \begin{equation}
    \IC(s, r) = \frac{\sum_{t=1}^n (R_t - \bar{R})(Q_t - \bar{Q})}
                     {\sqrt{\sum_{t=1}^n (R_t - \bar{R})^2}
                      \sqrt{\sum_{t=1}^n (Q_t - \bar{Q})^2}}
    \label{eq:spearman}
  \end{equation}
  For continuous distributions with no ties, $\bar{R} = \bar{Q} = (n+1)/2$
  and $\sum (R_t - \bar{R})^2 = n(n^2-1)/12$, giving the familiar formula:
  \[
    \IC(s, r) = 1 - \frac{6\sum_{t=1}^n d_t^2}{n(n^2 - 1)},
    \qquad d_t = R_t - Q_t
  \]
\end{definition}

\paragraph{Why Spearman over Pearson?}
The Pearson correlation $\Cor(s, r)$ is the natural measure of
\emph{linear} predictability.  However:
\begin{enumerate}[noitemsep]
  \item Log-return distributions are leptokurtic; a few large-return days
        dominate the Pearson correlation.  A single outlier can flip the sign.
  \item The relationship between a signal and $r_t^{(24h)}$ need not be
        linear.  The Spearman IC is invariant to any monotone transformation
        of either variable.
  \item Trading strategies based on signal \emph{rank} (enter when the signal
        is in the top quintile) are naturally assessed by rank correlation,
        not level correlation.
\end{enumerate}

\begin{remark}
  A negative IC is equally informative as a positive IC of the same
  magnitude: the signal should simply be traded on the opposite side.
  We report $|\IC|$ and the sign separately throughout.
\end{remark}

\subsection{The Breakeven Information Coefficient}
\label{sec:ICstar}

A statistically significant IC does not imply profitability.  The signal must
generate sufficient expected gross profit to cover round-trip transaction costs.
We now derive the minimum IC required.

\begin{proposition}[Breakeven IC]
  \label{prop:ICstar}
  Consider a trading strategy that enters long when a standardised signal
  $z_t \sim \mathcal{N}(0,1)$ is positive and short when $z_t$ is negative,
  exiting after $h$ bars.  Under the linear-IC model
  \begin{equation}
    r_t^{(h)} = \IC \cdot \sigma_r \cdot z_t + \varepsilon_t, \qquad
    \varepsilon_t \perp z_t,\quad
    \Var(r_t^{(h)}) = \sigma_r^2
    \label{eq:linearIC}
  \end{equation}
  the expected gross profit per trade equals
  $\IC \cdot \sigma_r \cdot \sqrt{2/\pi}$.  The signal is profitable
  against a round-trip cost $c_{\mathrm{RT}}$ if and only if
  \begin{equation}
    |\IC| > \IC^* \;\coloneqq\; \frac{c_{\mathrm{RT}}}{\sigma_r \sqrt{2/\pi}}
    \label{eq:ICstar}
  \end{equation}
\end{proposition}

\begin{proof}
  Under model~\eqref{eq:linearIC}, since $z_t \perp \varepsilon_t$:
  \[
    \E[r_t^{(h)} \mid z_t] = \IC \cdot \sigma_r \cdot z_t
  \]
  For a long trade conditional on $z_t > 0$:
  \begin{align*}
    \E[r_t^{(h)} \mid z_t > 0]
    &= \IC \cdot \sigma_r \cdot \E[z_t \mid z_t > 0]
    = \IC \cdot \sigma_r \cdot \frac{\E[z_t \cdot \mathbf{1}_{z_t > 0}]}
                                     {P(z_t > 0)}
  \end{align*}
  For $z_t \sim \mathcal{N}(0,1)$:
  \[
    \E[z_t \mid z_t > 0]
    = \frac{\int_0^\infty z\,\phi(z)\,\mathrm{d}z}{\tfrac{1}{2}}
    = 2\int_0^\infty z\,\phi(z)\,\mathrm{d}z
    = 2\left[-\phi(z)\right]_0^\infty
    = 2\phi(0)
    = \sqrt{\tfrac{2}{\pi}}
  \]
  By symmetry, a short trade ($z_t < 0$) earns $\E[-r_t^{(h)} \mid z_t < 0]
  = \IC \cdot \sigma_r \cdot \sqrt{2/\pi}$ as well.

  Expected net profit per trade equals gross minus round-trip cost:
  \[
    \E[\text{net profit}] = \IC \cdot \sigma_r \sqrt{2/\pi} - c_{\mathrm{RT}}
  \]
  Setting this to zero and solving for IC gives~\eqref{eq:ICstar}.\qed
\end{proof}

\begin{remark}[Threshold entry]
  If the strategy only trades when $|z_t| > c$ for some $c > 0$, the expected
  gross profit per trade is $\IC \cdot \sigma_r \cdot \E[z \mid z > c]$.  Since
  $\E[z \mid z > c] = \phi(c)/[1-\Phi(c)] > \sqrt{2/\pi}$ for all $c > 0$,
  the breakeven IC is \emph{lower} than~\eqref{eq:ICstar}.  The formula as
  stated is conservative: it applies to the strategy that enters on every bar.
\end{remark}

\begin{remark}[Pearson vs.\ Spearman]
  Proposition~\ref{prop:ICstar} is stated for the Pearson IC.  The Spearman
  IC satisfies $\IC_S \approx (6/\pi)\arcsin(\IC_P / 2)$, so
  $\IC_S \approx \IC_P$ for $|\IC| \ll 1$.  Since all empirical ICs in this
  report are in $[-0.20, +0.20]$, the same $\IC^*$ formula is applied.
\end{remark}

For the parameters used throughout Phase 5:
\begin{itemize}[noitemsep]
  \item $c_{\mathrm{RT}} = 6\bp$ (maker: $3\bp$ entry, $3\bp$ exit).
  \item $\sigma_r^{(24h)} \approx 2.135\%$ (empirical OOS volatility).
  \item $\IC^*_{\mathrm{maker}} \approx 0.035$.
\end{itemize}

The IC/IC* ratio gives the ``multiple of breakeven'': a ratio of $3\times$
means the signal generates three times the expected gross profit needed
to cover costs.

\subsection{The Overlapping-Return Inflation Problem}
\label{sec:overlap}

The 1-minute BTCUSDT dataset has $T \approx 527{,}040$ observations in the OOS
2024 window.  If we compute $r_t^{(24\text{h})} = \log(C_{t+1440}/C_t)$
for every $t$ and treat all pairs as independent, we overstate statistical
significance severely.

\paragraph{The autocorrelation structure of overlapping returns.}
Under the assumption that 1-minute log-returns are i.i.d.\ with
variance $\delta^2$, the 24h log-return decomposes as:
\[
  r_t^{(24\text{h})} = \sum_{j=0}^{1439} \Delta_j(t),
  \qquad \Delta_j(t) = \log\frac{C_{t+j+1}}{C_{t+j}}
\]
Two adjacent observations share all but one minute of their window.
More generally, for lag $k \leq 1440$:
\begin{align*}
  \Cov\!\left(r_t^{(24\text{h})}, r_{t+k}^{(24\text{h})}\right)
  &= \Cov\!\left(\sum_{j=0}^{1439}\Delta_j(t),\,
                 \sum_{j=0}^{1439}\Delta_j(t+k)\right)
   = (1440 - k)\,\delta^2
\end{align*}
since the two sums share exactly $1440 - k$ common terms (for $k < 1440$).
The autocorrelation at lag $k$ is:
\begin{equation}
  \Cor\!\left(r_t^{(24\text{h})}, r_{t+k}^{(24\text{h})}\right)
  = \frac{(1440 - k)\,\delta^2}{1440\,\delta^2}
  = 1 - \frac{k}{1440}, \qquad 0 \leq k \leq 1440
  \label{eq:overlap_autocor}
\end{equation}
This is a tent function that linearly decays from 1 at lag 0 to 0 at lag 1440.

\paragraph{Effective sample size.}
For a stationary series with autocorrelation function $\rho_k$, the effective
sample size for variance estimation is approximately:
\[
  n_{\mathrm{eff}} \approx \frac{T}{1 + 2\sum_{k=1}^{\infty} \rho_k}
\]
Substituting~\eqref{eq:overlap_autocor}:
\[
  1 + 2\sum_{k=1}^{1440-1}\!\!\left(1 - \frac{k}{1440}\right)
  = 1 + 2 \cdot \frac{1440 \cdot 1439}{2 \cdot 1440}
  = 1 + 1439
  = 1440
\]
Therefore:
\begin{equation}
  n_{\mathrm{eff}} \approx \frac{T}{1440} \approx 366
  \label{eq:neff}
\end{equation}
The effective number of independent 24-hour return observations is
approximately 366 per year --- one per day --- regardless of whether we
compute the return at every 1-minute bar or sample once per day.

\paragraph{Practical implication.}
A $t$-statistic computed from the full $n = 527{,}040$ series is inflated
by a factor of $\sqrt{1440} \approx 38$ relative to the correct statistic
based on $n_{\mathrm{eff}} \approx 366$.  \textbf{Solution:} sample every
$1{,}440$ bars ($n \approx 365$ non-overlapping daily observations per OOS
year) and apply the block bootstrap.

\subsection{The Block Bootstrap}
\label{sec:bootstrap}

The block bootstrap (Künsch, 1989; Liu \& Singh, 1992) provides valid inference
for statistics of dependent time series without parametric assumptions about
the dependence structure.

\begin{definition}[Non-Overlapping Block Bootstrap]
  Given a time series $(X_1, \ldots, X_n)$ and a statistic $T_n$:
  \begin{enumerate}
    \item Choose block length $\ell$ and form $m = \lfloor n/\ell \rfloor$
          non-overlapping blocks $B_j = (X_{(j-1)\ell+1}, \ldots, X_{j\ell})$
          for $j = 1, \ldots, m$.
    \item Draw $m$ blocks $B_{j_1}^*, \ldots, B_{j_m}^*$ uniformly at random
          with replacement from $\{B_1, \ldots, B_m\}$.
    \item Concatenate to form the bootstrap resample
          $X^* = (B_{j_1}^*, \ldots, B_{j_m}^*)$; trim to length $n$.
    \item Compute the bootstrap statistic $T_n^* = T_n(X^*)$.
    \item Repeat $B$ times.  The empirical distribution of
          $\{T_n^{*(b)}\}_{b=1}^B$ approximates the sampling distribution
          of $T_n$.
  \end{enumerate}
\end{definition}

For inference on $\IC$, the block bootstrap preserves temporal autocorrelation
within blocks (capturing weekly seasonality, funding-cycle effects, and
momentum effects) while permuting the dependence structure across blocks.

\paragraph{Block length selection.}
The block length $\ell$ should exceed the effective autocorrelation decay
length of the signal-return interaction.  For the daily non-overlapping series
used here, a block of $\ell = 21$ days (three calendar weeks) captures
relevant dependence: weekly trading patterns, bi-weekly options gamma pinning
cycles near expiry, and multi-week momentum effects.

\paragraph{$p$-value computation.}
For a one-sided test of $H_0 : \IC = 0$ against $H_1 : \IC > 0$:
\begin{equation}
  p = \frac{1}{B}\sum_{b=1}^B \mathbf{1}\!\left[\IC^{*(b)} \geq \widehat{\IC}\right]
  \label{eq:pvalue}
\end{equation}
where $B = 2{,}000$ resamples are used throughout.

\subsection{Pass Criteria (Pre-Committed)}

\textbf{IC screen:}
\begin{itemize}[noitemsep]
  \item Unfiltered: IC/IC* $> 0.5$ and correct sign.
  \item DVOL-filtered: ratio $> 1.0$ and bootstrap $p \leq 0.05$.
  \item Sub-period stability: IC same-sign in $\geq 3$ of 5 OOS half-years.
\end{itemize}

\textbf{Deep dive (additional gates for IC screen advances):}
\begin{itemize}[noitemsep]
  \item 95\% block-bootstrap CI lower bound strictly positive.
  \item Incremental IC (after N3$z$ control): $p \leq 0.05$.
  \item Standalone trade frequency $\geq$ N3's $1.2$ trades/week (Phase 5
        motivation gate: the signal must address the frequency problem).
\end{itemize}

\subsection{Data}

\begin{center}
\begin{tabular}{lll}
\toprule
Dataset & Source & Coverage \\
\midrule
1-minute BTCUSDT klines & Binance FAPI & 2023-01-01 to 2026-05-13 \\
BTC DVOL 1h & Deribit API & 2022-01-01 to 2026-05-15 \\
Bybit BTCUSDT funding (8h) & Bybit v5 API & 2022-01-01 to 2026-05-15 \\
OKX BTC-USDT-SWAP funding (8h) & OKX v5 API & 2026-02-15 to 2026-05-15 \\
\bottomrule
\end{tabular}
\end{center}

Bybit funding history: downloaded via \texttt{/v5/market/funding/history}
in 60-day windows (200 records per call), yielding 4{,}786 settlement
records.  OKX funding history: \texttt{/api/v5/public/funding-rate-history}
via cursor-based backward pagination, yielding only 276 records
($\approx 90$ days of retention --- an API limitation).

%% ============================================================
\section{S1: Realized Skewness}
\label{sec:s1}
%% ============================================================

\subsection{Signal Construction}

\begin{definition}[5-Day Realized Skewness]
  Let $r_t = \log(C_t / C_{t-1})$ be the 1-minute log return.
  The rolling 5-day ($W = 5 \times 1{,}440 = 7{,}200$ bars) realized
  skewness is estimated using the ratio of central moments:
  \begin{align}
    \hat{\mu}_t &= \frac{1}{W}\sum_{s=0}^{W-1} r_{t-s} \\
    \hat{\sigma}^2_t &= \frac{1}{W}\sum_{s=0}^{W-1}(r_{t-s} - \hat{\mu}_t)^2 \\
    \hat{\mu}^{(3)}_t &= \frac{1}{W}\sum_{s=0}^{W-1}(r_{t-s} - \hat{\mu}_t)^3 \\
    \hat{\gamma}_t &= \frac{\hat{\mu}^{(3)}_t}{(\hat{\sigma}^2_t)^{3/2} + \varepsilon}
    \label{eq:realskew}
  \end{align}
  where $\varepsilon = 10^{-20}$ prevents division by zero on
  zero-variance windows.  In vectorised implementation, the rolling
  central moments are computed as running sums to avoid quadratic
  time complexity.
\end{definition}

\begin{definition}[Skewness Signals]
  \begin{align}
    \text{S1a}_t &= -\hat{\gamma}_t
      \quad\text{(contrarian: negative skew $\to$ long)} \\
    \text{S1b}_t &= -\Delta\hat{\gamma}_t = \hat{\gamma}_{t-1440} - \hat{\gamma}_t
      \quad\text{(velocity: skew becoming more negative $\to$ long)}
  \end{align}
  S1a is positive (suggesting a long) when 5-day realized skewness is
  negative (many small gains, occasional large losses).  S1b is positive
  when skewness is becoming more negative (the fat left tail is
  growing in the recent window).
\end{definition}

\paragraph{Window length choice.}
The 5-day window ($W = 7{,}200$ bars) is chosen to capture the timescale
of individual liquidation cascade events.  A cascade typically plays out
over 12--48 hours; the 5-day window ensures the skewness estimate is
dominated by the most recent event without being swamped by older
background dynamics.  Longer windows (30-day) would dilute the signal;
shorter windows (1-day) would be too noisy at the individual bar level.

\paragraph{Descriptive statistics.}
Over the OOS daily frame:
S1a mean $= -0.021$, std $= 0.284$, range $[-1.39, +1.31]$.
$\Cor(\text{S1a}, N3z) = +0.126$ (low collinearity).

\subsection{IC Screen Results}

\begin{center}
\begin{tabular}{lrrrrrrl}
\toprule
Signal & IC & Ratio & $p$ & IC (filt.) & Ratio (filt.) & Stab. & Verdict \\
\midrule
\textbf{S1a level}  & $+0.058$ & $1.94\times$ & $0.039$ & $+0.097$ & $3.79\times$ & $4/5$ & \textbf{ADVANCE} \\
S1b velocity        & $-0.019$ & $0.63\times$ & $0.754$ & $-0.022$ & $0.86\times$ & $2/5$ & KILL \\
\bottomrule
\end{tabular}
\end{center}

\paragraph{S1b kill rationale.}
The velocity of skewness reversion (S1b) is in the wrong direction
($\IC = -0.019$) and fails $p = 0.754$.  The speed at which skewness
is becoming more negative adds no directional information about the 24h
return.  This is consistent with the hypothesis that the level of
skewness (how negative it is) matters for the exhaustion call, not its
rate of change (how fast it is getting more negative).

\paragraph{S1a advancing.}
S1a passes the unfiltered screen: IC $= +0.058$ ($1.94\times$),
$p = 0.039$ (below $0.05$), $4/5$ sub-period stability.  The DVOL
$\geq 50$ filter raises IC substantially to $+0.097$ ($3.79\times$),
consistent with the pattern from N3, P3, and Q5c: the signal is
concentrated in the fear regime.  S1a advances to the deep dive.

\subsection{S1a Deep Dive}
\label{sec:s1a_deep}

\begin{center}
\begin{tabular}{ll}
\toprule
Metric & Value \\
\midrule
OOS IC (non-overlapping daily) & $+0.058$ ($1.94\times$ ratio) \\
Block-bootstrap $p$ & $0.039$ \\
\textbf{95\% CI lower bound} & $\mathbf{-0.008}$ \textbf{--- fails: not strictly positive} \\
DVOL $\geq 50$ filtered IC & $+0.097$ ($3.79\times$), $p = 0.021$ \\
Sub-period stability & $4/5$ OOS half-years \\
Incremental IC $p$ (after N3$z$) & $0.234$ --- not significant \\
$\Cor(\text{S1a}, N3z)$ & $+0.126$ \\
\midrule
Standalone Sharpe (S1a $\geq 0.5$, DVOL $\geq 50$) & $+1.55$ ($n = 80$ trades, $0.7$/wk) \\
Standalone Sharpe \emph{excl.\ N3$z$ regime} & $-2.26$ --- loses money \\
\midrule
Combined (S1a $\geq 0.5$ AND N3$z \geq 0$, DVOL $\geq 50$) & $+8.85$ ($n = 35$ trades, $0.3$/wk) \\
\bottomrule
\end{tabular}
\end{center}

\subsubsection{Gate-by-Gate Assessment}

\textbf{Gate 1 --- 95\% CI lower bound.}
The block-bootstrap 95\% confidence interval for the IC is $(-0.008, +0.124)$.
The lower bound of $-0.008$ is \emph{negative}: the sampling distribution
of the IC extends below zero.  We cannot rule out that the true IC is
zero or negative.  This gate fails.

To understand why the $p$-value ($0.039$) and the CI lower bound
($-0.008$) appear contradictory: the $p$-value tests the one-sided
null ($\IC \leq 0$), and $p = 0.039$ means we barely reject it.
The CI is a two-sided interval; its lower bound being $-0.008$ is
mathematically consistent with $p = 0.039$ and means the evidence
is marginal.  Both the $p$-value and the CI tell the same story:
the IC of $+0.058$ is at the boundary of statistical credibility
with $n \approx 365$ observations.

\textbf{Gate 2 --- Incremental IC after N3$z$ control.}
Projecting S1a onto N3$z$ and taking the residual:
\[
  \tilde{\text{S1a}}_t = \text{S1a}_t - \hat{\beta}_0 - \hat{\beta}_1 N3z_t
\]
The IC of $\tilde{\text{S1a}}$ against 24h returns has $p = 0.234$.
After removing the N3$z$ component, S1a has no significant
residual predictive power.  This is the critical finding: \textbf{the
IC of S1a is driven by its correlation with N3$z$}, not by the skewness
mechanism independently.

Why does this happen?  Both S1a and N3$z$ are elevated on high-DVOL
days.  The correlation is $+0.126$ (small but non-negligible).  Days of
negative realised skewness tend to coincide with high DVOL because both
are elevated by the same liquidation cascade events.  After controlling
for N3$z$, there is no incremental directional information in the
skewness beyond what is already captured by the fear-resolution mechanism.

\textbf{Gate 3 --- Standalone performance outside the N3 regime.}
When N3$z$ is not firing (below the activation threshold), S1a alone
achieves Sharpe $= -2.26$.  This means \emph{S1a loses money on days
when N3 is not active}.  On these days, negative realised skewness
reflects momentum continuation (the left-skewed return distribution
persists as the downtrend continues) rather than the exhaustion-bounce
mechanism that makes the contrarian bet profitable.

\textbf{Gate 4 --- Phase 5 frequency gate.}
Even the combined rule (S1a AND N3$z$, DVOL $\geq 50$) fires only
$0.3$ times per week --- substantially less than N3 alone ($1.2$/wk).
Phase 5 was motivated by the desire to \emph{increase} trade frequency;
S1a decreases it.  Adding S1a as a filter on N3 improves quality
(Sharpe $+8.85$ vs N3's $+2.95$) at the cost of frequency (0.3 vs 1.2
trades/week), the opposite of what Phase 5 sought.

\subsection{S1a Kill Rationale and Residual Value}

\textbf{S1a is formally killed as a standalone signal.}

The four deep-dive gates tell a consistent story:
\begin{itemize}
  \item The statistical evidence is marginal (CI crosses zero).
  \item The predictive power is not independent of N3$z$.
  \item Outside the N3 regime, S1a loses money.
  \item S1a reduces trade frequency rather than increasing it.
\end{itemize}

However, the combination N3$z \geq 0$ AND S1a $\geq 0.5$ (DVOL $\geq 50$)
achieves Sharpe $+8.85$ on 35 trades.  This is noted: S1a may be useful
as a \emph{quality filter} on N3 entries in a future phase if the
frequency constraint is relaxed.  But it cannot be the standalone
higher-frequency signal Phase 5 sought.

%% ============================================================
\section{S2: Pre-Settlement Flow}
\label{sec:s2}
%% ============================================================

\subsection{Signal Construction}

\begin{definition}[Taker Order Flow Imbalance]
  Let $V^{\mathrm{buy}}_t$ and $V_t$ be the 1-minute taker-buy volume
  and total volume from Binance klines.  The taker fraction and signed
  flow at bar $t$ are:
  \begin{align}
    \phi_t &= \frac{V^{\mathrm{buy}}_t}{V_t \vee \varepsilon}
    \quad\in [0,1] \\
    f_t &= 2\phi_t - 1 \quad\in [-1, +1]
  \end{align}
  where $\varepsilon = 10^{-8}$ prevents division by zero on zero-volume bars.
  $f_t = +1$ means 100\% of volume was taker-buy; $f_t = -1$ means 100\%
  was taker-sell.
\end{definition}

\begin{definition}[Pre-Settlement Flow Signals]
  The 30-minute cumulative signed flow before each settlement is:
  \begin{equation}
    F^{(30)}_t = \sum_{s=t-29}^{t} f_s
    \quad\in [-30, +30]
  \end{equation}
  This is computed at all bars but evaluated only at the 480 bars
  ($= 8\text{h} \times 60$) preceding each settlement.
  Let $\lambda_t$ be the current 8-hour funding rate.  The signals:
  \begin{align}
    \text{S2a}_t &= \phantom{-}F^{(30)}_t \times (-\lambda_t)
      \quad\text{(momentum)} \\
    \text{S2b}_t &= -F^{(30)}_t \times (-\lambda_t)
      \quad\text{(contrarian)}
  \end{align}
  S2a is positive when taker selling coincides with a positive funding
  rate: sellers are anticipated informed traders.  S2b reverses the
  sign: the selling is mechanical arbitrage that reverts post-settlement.
\end{definition}

\subsection{IC Screen Results}

\begin{center}
\begin{tabular}{lrrrrrrl}
\toprule
Signal & IC & Ratio & $p$ & IC (filt.) & Ratio (filt.) & Stab. & Verdict \\
\midrule
S2a momentum   & $+0.001$ & $0.04\times$ & $0.484$ & $+0.006$ & $0.23\times$ & $3/5$ & KILL \\
S2b contrarian & $+0.014$ & $0.47\times$ & $0.153$ & $+0.002$ & $0.02\times$ & $2/5$ & KILL \\
\bottomrule
\end{tabular}
\end{center}

\subsection{Kill Rationale}

Both S2 variants are cleanly killed.  S2a's IC of $+0.001$ ($0.04\times$)
is effectively zero.  S2b's unfiltered ratio of $0.47\times$ is just below
the $0.5\times$ threshold, and the DVOL-filtered IC collapses to
$0.02\times$ --- statistically and economically indistinguishable from noise.

The empirical result confirms the theoretical expectation from
Section~\ref{sec:settlement_theory}: the pre-settlement mechanism
operates on 2--4 hour timescales and is invisible in the 24h return
target.  At the daily frequency, three daily settlement events are
aggregated into a single signal observation, averaging away any
directional content.

\begin{remark}
  Testing S2 at the 8-hour forward return horizon (sampling at each
  settlement event, $n \approx 3 \times 365 = 1095$ OOS observations/year)
  would be the correct evaluation framework.  This requires:
  (1) a settlement-event-level dataset, (2) 8h forward returns aligned
  to each settlement, and (3) a modified block bootstrap block length
  of $\approx 3 \times 21 = 63$ settlement events.  This is a Phase 6
  scope item.
\end{remark}

%% ============================================================
\section{S3: Cross-Exchange Funding Divergence}
\label{sec:s3}
%% ============================================================

\subsection{Signal Construction}

\begin{definition}[Funding Rate Divergence]
  For each 8-hour settlement bar, the raw Bybit--Binance divergence is:
  \begin{equation}
    \delta^{(\mathrm{Byb})}_t = F^{(\mathrm{Byb})}_t - F^{(\mathrm{Bin})}_t
    \label{eq:diverg}
  \end{equation}
  where $F^{(\mathrm{Byb})}_t$ and $F^{(\mathrm{Bin})}_t$ are the
  respective 8-hour funding rates, aligned to the nearest Binance
  settlement timestamp.

  The signal is $z$-scored over a 30-settlement (90-day) rolling window:
  \begin{equation}
    \text{S3a}_t = \frac{\delta^{(\mathrm{Byb})}_t - \mu_{30}(\delta^{(\mathrm{Byb})})}
                        {\sigma_{30}(\delta^{(\mathrm{Byb})})}
    \label{eq:s3a}
  \end{equation}
  S3b uses the OKX--Binance divergence with the same construction.
\end{definition}

\subsection{IC Screen Results}

\begin{center}
\begin{tabular}{lrrrrrrl}
\toprule
Signal & IC & Ratio & $p$ & IC (filt.) & Ratio (filt.) & Stab. & Verdict \\
\midrule
S3a Bybit--Binance & $-0.034$ & $1.14\times$ & $0.840$ & $-0.018$ & $0.71\times$ & $2/5$ & KILL \\
S3b OKX--Binance   & \multicolumn{6}{l}{\emph{Untestable: OKX API retains $\approx 90$ days only (7\% OOS coverage)}} & KILL \\
\bottomrule
\end{tabular}
\end{center}

\subsection{Kill Rationale}

\paragraph{S3a (Bybit--Binance).}
The IC is in the \emph{wrong} direction: $-0.034$, with $p = 0.840$
(one-sided against the positive direction).  A positive Bybit--Binance
divergence is weakly associated with \emph{negative} subsequent Binance
returns --- the opposite of the arbitrage-flow prediction.

The most likely explanation: at the 8-hour aggregation level, funding
rate divergences are highly noisy.  The Bybit and Binance rates are
computed from different underlying spot indices, different perpetual
order books, and different user bases.  By the time the 8-hour average
is reported, any cross-exchange arbitrage opportunity has already been
closed.  The residual divergence in the reported rates is dominated by
the difference in reference spot indices (not by speculative demand)
and contains no actionable directional information.

The $2/5$ sub-period stability confirms the IC sign is not consistent
across periods.

\paragraph{S3b (OKX--Binance).}
The OKX v5 public endpoint (\texttt{/api/v5/public/funding-rate-history})
returned 276 records covering approximately 90 days.  The OKX API does
not retain longer history on this endpoint.  Against the $\approx 3{,}800$-bar
OOS window, this is 7\% coverage --- insufficient for a meaningful IC
screen.  S3b is formally recorded as \textbf{untestable}.

%% ============================================================
\section{Synthesis: The High-Frequency Problem}
\label{sec:synthesis}
%% ============================================================

\subsection{Why All Phase 5 Signals Failed}

Phase 5 was motivated by a specific objective: find signals with higher
trade frequency than N3 ($1.2$ trades/week) and P3.  The fundamental
finding is that \emph{no such signal exists at the daily-resolution IC
framework used in this programme.}

Table~\ref{tab:kills} summarises all kills:

\begin{table}[H]
\centering
\caption{Phase 5 complete kill log}
\label{tab:kills}
\begin{tabular}{lll}
\toprule
Signal & Mechanism tested & Kill reason \\
\midrule
S1a Skewness level & Liquidation exhaustion & CI lower $= -0.008$; incremental IC $p = 0.234$ \\
S1b Skewness velocity & Cascade acceleration & Wrong direction, $p = 0.754$ \\
S2a Pre-settlement (momentum) & Informed pre-settlement selling & Near-zero IC ($0.04\times$) \\
S2b Pre-settlement (contrarian) & Mechanical arb reversal & Below threshold ($0.47\times$) \\
S3a Bybit--Binance & Cross-exchange arb flow & Wrong direction, $p = 0.840$ \\
S3b OKX--Binance & Cross-exchange arb flow & OKX API retains $\approx 90$ days only \\
\bottomrule
\end{tabular}
\end{table}

\subsection{The Structural Concentration of Edge in the Fear Regime}

The S1a deep dive reveals the structural reason why higher-frequency
signals are difficult to find within the current framework:

\begin{itemize}
  \item S1a has IC $= +0.058$ overall.
  \item S1a has IC $\approx 0$ (or negative) in the low-DVOL regime.
  \item S1a has IC $= +0.097$ in the DVOL $\geq 50$ regime.
  \item After controlling for the DVOL regime (via N3$z$), S1a's
        incremental IC is $p = 0.234$ --- not significant.
\end{itemize}

The inescapable conclusion is that \emph{every mechanism tested so far
that predicts positive 24h BTC returns at the daily frequency does so
because it is a proxy for the same underlying state: elevated fear
(high DVOL) followed by mean reversion of that fear.}

This is a structural feature of the market, not a measurement artifact.
The economics of perpetual futures mean that the primary source of
exploitable edge at the 24h daily-frequency horizon is the
liquidation-exhaustion cycle in high-DVOL environments.  N3 and P3
each measure different observable manifestations of this cycle
(IV $z$-score, and OI-price regime, respectively).  S1a (skewness)
is a third manifestation of the same cycle, which is why it has IC
after conditioning on DVOL but no incremental IC after conditioning on N3$z$.

\subsection{What Higher-Frequency Research Requires}

A genuine higher-frequency strategy in crypto perpetuals requires
a different research infrastructure:

\begin{enumerate}
  \item \textbf{Intraday forward returns.}  Evaluating S2 at the 8h
        horizon requires $\sim 1{,}000$ OOS observations per year at
        3 settlements/day, and a modified statistical framework
        (shorter blocks, settlement-aligned sampling).

  \item \textbf{Tick-level data for adverse selection.}  Pre-settlement
        flow signals require the ability to distinguish informed from
        uninformed taker flow, which requires bid-ask spread and depth
        data at the order book level --- not just the aggregated 1-minute
        kline data.

  \item \textbf{Explicit settlement mechanics modelling.}  The pre-settlement
        window begins 480 minutes before the settlement; exactly modelling
        which bars fall within this window requires precise settlement time
        data, not just the funding rate time series.

  \item \textbf{Alternative liquidation data.}  The Q6 (liquidation
        exhaustion) hypothesis (Phase 4) remains untested.  Intraday
        liquidation data from Coinalyze or CoinGlass would allow testing
        Q6 at the event level (1--4h after a large cascade).
\end{enumerate}

These infrastructure items define the Phase 6 research scope.

%% ============================================================
\section{Summary and Kill Log}
\label{sec:summary}
%% ============================================================

\begin{center}
\begin{tabular}{lll}
\toprule
Signal & Verdict & Primary reason \\
\midrule
S1a Realized skewness (contrarian) & Kill & CI lower $= -0.008$; incremental IC $p = 0.234$;
  loses money outside N3 \\
S1b Skewness velocity & Kill & Wrong direction ($2/5$ stability, $p = 0.754$) \\
S2a Pre-settlement flow (momentum) & Kill & Near-zero IC ($0.04\times$), $p = 0.484$ \\
S2b Pre-settlement flow (contrarian) & Kill & Below threshold ($0.47\times$); DVOL-filtered $0.02\times$ \\
S3a Bybit--Binance divergence & Kill & Wrong direction ($2/5$, $p = 0.840$) \\
S3b OKX--Binance divergence & Untestable & OKX API retains $\approx 90$ days only \\
\midrule
\multicolumn{3}{l}{\textbf{Phase 5 net result: zero new signals validated.  N3 + P3 stack unchanged.}} \\
\bottomrule
\end{tabular}
\end{center}

Research scripts: \texttt{data/download\_phase5.py},
\texttt{research/active/p5/27\_phase5\_ic\_screen.py},
\texttt{research/active/p5/28\_s1a\_deep\_dive.py}.

%% ============================================================
\section{Conclusion}
\label{sec:conclusion}
%% ============================================================

Phase 5 tested three signal families and validated none.  The key results:

\begin{enumerate}
  \item \textbf{Realized skewness (S1a) is a proxy for N3, not an
        independent signal.}  The 5-day negative realized skewness has IC
        $= +0.058$ at the IC screen, but fails the deep dive: the 95\%
        CI crosses zero, and after controlling for N3$z$, the incremental
        IC has $p = 0.234$.  S1a loses money ($\text{Sharpe} = -2.26$)
        on days when N3 is not active.  Every mechanism tested in Phase 5
        that predicts positive BTC returns at the daily frequency does so
        because it correlates with the same high-DVOL exhaustion state
        that N3 already captures.

  \item \textbf{Pre-settlement flow (S2) is invisible at daily resolution.}
        The settlement-arbitrage mechanism exists in theory, but the
        relevant forward return window (2--4h post-settlement) is
        subsumed by the 24h IC evaluation frame.  Near-zero ICs ($0.04\times$
        and $0.47\times$) confirm the signal is indistinguishable from
        noise at daily sampling.

  \item \textbf{Cross-exchange funding divergence (S3) has no edge at
        the 8-hour aggregation level.}  Well-capitalised arbitrageurs
        close Bybit--Binance funding gaps within minutes, leaving no
        signal in the reported 8-hour rates.  S3a's IC is in the wrong
        direction ($-0.034$, $p = 0.840$).

  \item \textbf{The high-frequency problem is structural, not accidental.}
        Edge in BTCUSDT perpetuals at the 24h daily-frequency horizon
        is concentrated in the high-DVOL fear regime.  Finding genuinely
        higher-frequency alpha requires intraday return targets, tick-level
        data, and explicit settlement-mechanics modelling.  This is Phase 6.
\end{enumerate}

The production signal stack remains N3 and P3.  S1a's combined performance
with N3 (Sharpe $+8.85$ on 35 trades, $0.3$/week) is noted as a
potential quality-filter application for N3 entries, but is not deployed
due to the frequency reduction and the statistically marginal IC evidence.

\vspace{1em}
\noindent\textit{Data: Binance FAPI 1m klines (2023--2026); Deribit BTC
DVOL 1h (2022--2026); Bybit BTCUSDT 8h funding rates (2022--2026,
4{,}786 records); OKX BTC-USDT-SWAP 8h funding rates ($\approx$90 days,
276 records).}

%% ============================================================
\section*{References}
%% ============================================================
\begingroup
\setlength{\parindent}{-2em}
\setlength{\leftskip}{2em}
\setlength{\parskip}{6pt}
\small

Amaya, D., Christoffersen, P., Jacobs, K., \& Vasquez, A. (2015). Does realized skewness predict the cross-section of equity returns? \textit{Journal of Financial Economics}, \textit{118}(1), 135--167. \url{https://doi.org/10.1016/j.jfineco.2015.02.012}

Harvey, C. R., \& Siddique, A. (2000). Conditional skewness in asset pricing tests. \textit{The Journal of Finance}, \textit{55}(3), 1263--1295. \url{https://doi.org/10.1111/0022-1082.00247}

Künsch, H. R. (1989). The jackknife and the bootstrap for general stationary observations. \textit{The Annals of Statistics}, \textit{17}(3), 1217--1241. \url{https://doi.org/10.1214/aos/1176347265}

Liu, R. Y., \& Singh, K. (1992). Moving blocks jackknife and bootstrap capture weak convergence. In R. LePage \& L. Billard (Eds.), \textit{Exploring the Limits of Bootstrap} (pp. 225--248). Wiley.

Makarov, I., \& Schoar, A. (2020). Trading and arbitrage in cryptocurrency markets. \textit{Journal of Financial Economics}, \textit{135}(2), 293--319. \url{https://doi.org/10.1016/j.jfineco.2019.07.001}

\endgroup

\end{document}
